\section{Measuring index calculus against rho}\label{sec:method}

The programme's first lesson was that index-calculus numbers from
different studies did not compose: one study timed a decomposition, another
a relation, a third a whole multi-target batch with shared tables, each on
a different host and in a different unit. Four mechanisms, built in this
order, fixed that.

\subsection{The boundary ledger (10 September)}

\art{docs/ic/BOUNDARY_TARGETS.md} and its machine twin
\art{docs/ic/boundary_targets.json} give each regime (Koblitz, binary,
prime) and each pipeline stage (\texttt{factor\_base},
\texttt{decomposition}, \texttt{relation\_yield}, \texttt{rank},
\texttt{end\_to\_end\_dlp}, \texttt{vs\_rho}) a current record, a next
target and acceptance gates. A row promotes only when its evidence JSON
carries every field the measurement schema requires for that stage---for
example, a \texttt{decomposition} claim that beats an algebraic frontier
must report the first-fall degree, and a \texttt{vs\_rho} claim needs one
target, the identical public point for both methods, verified scalars,
five exclusive IC online phases, the automorphism discount, matched
resources and a host-isolation receipt. Missing fields fail closed. The
primary \texttt{vs\_rho} question is the online cost of \emph{one}
previously unseen target after reusable IC preparation; multi-target and
amortised results are secondary and must be declared as such.

\subsection{ecbench (from 2 October)}\label{sec:ecbench}

\emph{ecbench} (\art{src/cryptanalysis/ecbench/}, \art{src/bin/ecbench.rs},
documented in \art{docs/ecbench/README.md}) runs every ECDLP method behind
one interface, on the same single-target workloads, in one unit, into
sealed records.

\paragraph{Unit.}
\[
  S \;=\; \frac{\text{total group-addition equivalents (GAE)}}{\sqrt r},
  \qquad
  S_{\text{floor}}=\sqrt{\pi/(2A)}.
\]
An addition or doubling is one GAE; a scalar multiplication is charged as
the additions it performs; native work (lookups, row operations, square
roots, Artin--Schreier solves, word XORs) is priced at a pinned conversion
where the repository has one and is otherwise recorded as
\emph{uncharged}, which marks the record's $S$ as a \emph{lower bound}.
Plain rho has $S\approx1.25$.

\paragraph{What a measurement is.} One workload is one curve, subgroup,
generator and target; nothing survives between executions, so every table
a method builds is inside its cost (\emph{cold}). Beside the cold $S$
every record carries the one-target \emph{online window}, split for index
calculus into the five exclusive phases \texttt{target\_query},
\texttt{target\_PDP}, \texttt{target\_relation\_check},
\texttt{target\_descent}, \texttt{target\_recovery\_check}. Targets are
either planted or \emph{public}---hashed to the subgroup so that nobody
knows the logarithm---and the runner verifies $[k]G=Q$ in its own process;
\texttt{ecbench verify} checks it again from the files. Failed, exhausted
and timed-out runs are rows, never dropped.

\paragraph{Identity, sealing and replay.} Every object (curve, factor
base, method parameters, spec, session, host environment class) is named by
a hash of what it is (\texttt{ECBS1h\ldots} for sessions,
\texttt{ECBND1h\ldots} for bounds). Each execution appends one sealed
record; \texttt{ecbench verify} replays measured records and emits an audit
receipt whose SHA-256 is the replay certificate. Operation counts are
deterministic, so a replay must match to the last digit; a mismatch is
preserved as a failed audit (\S\ref{sec:corrections}).

\paragraph{Isolation.} A run earns L0 (recorded under a host capsule), L1
(pinned, affinity read back, no foreign SMT sibling, CPU/wall $\le1.05$),
L2 (quiet: zero steal, foreign busy time $\le1\%$, run-queue delay
$\le0.5\%$, $\le50$ preemptions/s) or L3 (isolated cores, fixed governor,
bare metal). Wall time is admitted only at the spec's required level
(default L2) with at least five admitted pairs and an A/A interval;
otherwise it is printed as \emph{descriptive}. Operation counts need no
isolation. macOS hosts earn L0; a GitHub-hosted runner earned L1 on 72/72
runs and matched a macOS session's operation counts to the last digit.
\emph{isolab} prepares lab hosts for L3 and cross-host replay.

\paragraph{Statistics.} The comparable statistic is the mean $S$ (the
floor is an expectation). Comparisons are paired by (workload, round) with
shared seeds, reported as a ratio of totals $\sum S_B/\sum S_A$, never a
mean of ratios, with a two-stage cluster bootstrap that resamples
workloads, then pairs. An operations A/A is exactly 1 by construction and
checks determinism. A claim about scaling fits exponents over at least
four sizes.

\paragraph{Bounds and the frontier.} \texttt{ecbench} also records sealed
\emph{bounds}---constant and exponent with intervals, per phase---and a
challenge protocol whose verdict says whether a candidate advances,
trades, matches or regresses an incumbent and at which level (exponent,
constant, primitive). The incumbent for cold one-target Koblitz operations
is \texttt{rho.signed\_frobenius\_strong} (bound
\texttt{ECBND1h05a2ea918f3c}): $1.089\times$ the floor $[0.927,1.261]$,
fitted exponent $\alpha=0.451$ $[0.399,0.512]$, 144 runs at six sizes
(\art{research/ecbench_frontier_20261013/}).

\begin{definition}[comparability]\label{def:domain}
A performance point is keyed by the tuple
\begin{align*}
D=\bigl(&\text{curve and subgroup},\ \text{target law},\
  \text{cold}\mid\text{online}\mid\text{stage},\ \text{unit and unpriced work},\\
  &\text{set-up amortisation }q,\ \text{memory envelope},\
  \text{correctness/replay/isolation class}\bigr).
\end{align*}
Two points are ratio-comparable only when every component of $D$ except
the method matches; a method is on the frontier of a domain when no admitted point is
no worse on every axis and strictly better on one.
\end{definition}
\noindent This rules out the four collapses we saw most often: online wall
as cold cost; word XORs, probes, field operations and GAE as
interchangeable; a lower bound as a complete cost; an L0 wall ratio as an
admitted speed-up.

\subsection{The pricing band}\label{sec:band}

Almost no cost in a record is complete, and the rho baselines are no
exception. Rho, BSGS and the kangaroo leave canonicalisations, partition
hashes and table inserts and lookups unpriced. An index-calculus arm is
complete only on curves that have a pinned timing calibration
(\art{docs/ic/calibration.json}), which converts lookups, row operations,
Artin--Schreier solves, Frobenius maps and solver word XORs to GAE at
measured ratios of roughly $10^{-3}$ to $10^{-1}$. A ``complete'' index
calculus set against a lower-bound rho is therefore lopsided, and several
earlier comparisons in this repository were of that kind.

Since 15 October every cross-method comparison uses a \emph{pricing band}
(repository rule, \texttt{AGENTS.md} \S5a; native implementation
\texttt{ecbench band}). For a record with charged cost $G$, calibration-priced
native cost $C$ over $n_C$ operations, and $U$ operations behind its
unpriced counters,
\[
  S(w)=\frac{G-C+w\,(U+n_C)}{\sqrt r},\qquad w\in\{0,\ 0.1,\ 1\}
  \text{ GAE per native operation},
\]
applied alike to every arm, references included. Two arms are ranked only
if their order is the same at every weight; otherwise the pair is reported
as unresolved. An arm that names an unpriced counter without counting it is
a lower bound and cannot be ranked. The band brackets the pinned
calibrations from below and above, so a ranking that survives it does not
depend on the host the calibration came from.

\subsection{Claim hygiene and change classes}

Every study directory states a protocol before it runs (frozen inputs are
committed first), keeps failures, and closes with a decision in one of the
four classes of \S\ref{sec:intro}. A controlled-experiment rule requires a
negative control from the right null population, a positive control that
proves the detector can see the effect, and the same instrument on every
arm before any comparative statement. Blockers are recorded as named
problems (\art{problems/}) with exit criteria rather than silently dropped.
Measurement code is native (Rust); Python only lays out recorded values.

\subsection{Corrections the accounting forced}\label{sec:corrections}

Building the harness and the ledger exposed errors in earlier claims,
including our own. We list them because they change how earlier numbers
must be read.
\begin{enumerate}
\item \emph{Online Koblitz ratios relabelled} (7--8 October,
  \art{docs/ic/PLAN_IC_ACCOUNTING_FIXES_20261007.md}). Same-point online
  ratios of $124.85\times$ ($n=61$), $2{,}864.1\times$ ($n=71$),
  $1{,}189.7\times$ ($n=73$) and $8.2\times$ ($n=83$) had verified answers
  but no host-isolation receipt; they are \emph{exploratory}. Only the
  later 64-target pinned panel at $n=61$ is admitted (\S\ref{sec:koblitz}).
\item \emph{ECC2K-130 extrapolation constants corrected}
  (\art{RESEARCH_ECC2K130_IC_FEASIBILITY.md}): the structural model
  predicts $C=0.25$ in $\text{probes/relation}\approx C\,r/(n^2K^2)$; the
  fitted $C$ is $0.27$ at $n=61$ but $1.2$--$1.5$ at $n=71,73$, an
  unexplained gap now stated rather than absorbed.
\item \emph{``F4'' that was Buchberger}
  (\art{research/ecbench_rootfinding_assessment_20261012/REPORT.md}): the
  NIST benchmark of the root-finding manuscript printed its baseline as F4
  but called \texttt{groebner\_f4::buchberger}; its $3{,}021$--$5{,}802\times$
  are ratios to Buchberger. The manuscript's claim that no implementation
  can use fewer than linear field operations per ideal is too broad once a
  reusable pair table is allowed.
\item \emph{Unpriced word XORs} (\art{research/ecbench_rootfind_faster_20261014/}).
  Inherited F4's Boolean word XORs were counted but not charged, so its
  $S=16.7$ at Koblitz $n=13$ was a lower bound that looked competitive
  with rho ($14.4$). Pricing them gives $S=421.4$.
\item \emph{Cross-host counter drift}
  (\art{research/ecbench_frontier_20261013/audits/}): a replay-all audit
  of the root-finder session on x86\_64 matched 120/168 records; all 48
  inherited-F4 rows differed in solver word-XOR counts (e.g.\ 301{,}961
  versus 300{,}531), with every scalar, trial count and charged GAE
  unchanged. The failed receipt is kept; the F4 stage comparison is
  historical until the drift is explained.
\item \emph{A soundness bug} (\art{docs/ic/GROEBNER_F4F5F6_AUDIT_20261010.md}
  \S2.1): a degree-bounded \texttt{f4\_fp::solve} could certify solutions
  of an inconsistent system (over $\FF_5$, $\{x^3+1,xy-1,x+y\}$ returned
  two ``solutions''). Fixed by evaluating every candidate on the input.
\item \emph{F6-IC charged cost} (\S\ref{sec:groebner}): a stage win in
  word XORs became a $64.6$--$1{,}353\times$ loss in charged $S$ once its
  batched geometry was counted---recorded as an accounting result, not an
  advance.
\item \emph{Lopsided pricing} (\S\ref{sec:band}): comparisons of a
  calibration-priced index calculus with a lower-bound rho. On the
  frontier campaign's train panel at Koblitz $n=13$ the compact-orbit arm,
  earlier quoted at $1.08\times$ rho as a lower bound, is $0.93\times$,
  $1.53\times$ and $5.7\times$ the strong rho at $w=0$, $0.1$ and $1$
  (columns 2, four train targets,
  \art{research/ic_frontier_campaign_20261015/select/SELECTION.md}): below
  rho with native work free, above it otherwise, so the pair is unresolved
  inside the band rather than a win for either side.
\end{enumerate}
